\begin{lemma}[Exact hidden-tag usage law]\label{lem:step-009-usage-law}
Under the ideal hidden-coordinate tag interface, for every pair of
integers \(k\ge2,n\ge1\), the count \(U\) in (2) satisfies, for each
\(j\in[k]\) and each \(u\in\{0,\ldots,n\}\),

\[
\Pr\{U=u\mid J=j\}
=\binom nu\left(\frac1k\right)^u
 \left(1-\frac1k\right)^{n-u}.
\tag{4}
\]

Consequently \(U\sim\operatorname{Bin}(n,1/k)\) unconditionally. The law
includes \(U=0\), and its trial count is exactly \(n\).
\end{lemma}

\begin{proof}
Fix \(j\in[k]\). By the tag interface, conditional on \(J=j\),
the variables

\[
B_\ell:=\mathbf1\{I_\ell=j\},\qquad \ell\in[n],
\]

are mutually independent, and

\[
\Pr\{B_\ell=1\mid J=j\}=\Pr\{I_\ell=j\mid J=j\}=\frac1k.
\]

Thus \(U=\sum_{\ell=1}^nB_\ell\). For any fixed set of exactly \(u\)
indices, the probability that precisely those indicators equal one is
\(k^{-u}(1-1/k)^{n-u}\). The \(\binom nu\) choices are disjoint, proving
(4). The right-hand side does not depend on \(j\). Since
\(\Pr(J=j)=1/k\),

\[
\begin{aligned}
\Pr\{U=u\}
&=\sum_{j=1}^k\Pr\{J=j\}\Pr\{U=u\mid J=j\}\\
&=\binom nu\left(\frac1k\right)^u
 \left(1-\frac1k\right)^{n-u}.
\end{aligned}
\]

This is the asserted binomial mass function. In particular, no random block
instance, feature, label, learner output, or stopping rule changes the law.
\end{proof}

\begin{lemma}[Falling-factorial binomial tail bound]\label{lem:step-009-factorial-tail}
Let \(V\sim\operatorname{Bin}(n,p)\), where \(n\ge1\) and
\(p\in[0,1]\). For every integer \(v\ge0\), set \((v)_0:=1\); for every
integer \(r\ge1\), write

\[
(v)_r:=v(v-1)\cdots(v-r+1),
\]

with \((v)_r=0\) when \(r>v\). Then, for every integer \(r\ge1\),

\[
\mathbb E(V)_r=(n)_r p^r
\tag{5}
\]

and

\[
\Pr\{V\ge r\}
\le \frac{(n)_r p^r}{r!}
\le \frac{(np)^r}{r!}.
\tag{6}
\]

Equations (5)--(6) remain valid when \(r>n\), in which case both the event
and \((n)_r\) vanish.
\end{lemma}

\begin{proof}
Represent \(V=\sum_{\ell=1}^n Z_\ell\) for independent
Bernoulli\((p)\) variables. For every outcome with \(V=v\), the number of
ordered \(r\)-tuples of pairwise distinct successful trial indices is
exactly \((v)_r\). Hence the pointwise identity

\[
(V)_r
=\sum_{\substack{(\ell_1,\ldots,\ell_r)\in[n]^r\\
                    \ell_1,\ldots,\ell_r\ \mathrm{pairwise\ distinct}}}
  \prod_{s=1}^r Z_{\ell_s}
\tag{7}
\]

holds, with an empty sum if \(r>n\). There are \((n)_r\) ordered tuples,
and independence gives expectation \(p^r\) for every product in (7),
proving (5).

If \(V\ge r\), then

\[
(V)_r=r!\binom Vr\ge r!,
\]

whereas if \(V<r\), both \(\mathbf1\{V\ge r\}\) and \((V)_r\) are zero.
Therefore

\[
\mathbf1\{V\ge r\}\le\frac{(V)_r}{r!}
\]

pointwise. Taking expectations and using (5) gives the first inequality in
(6). Finally, each of the \(r\) nonnegative factors of \((n)_r\) is at
most \(n\) when \(r\le n\), so \((n)_r\le n^r\). When \(r>n\), the
left side is zero. This proves the second inequality in every case.
\end{proof}

\begin{lemma}[Floor-eight small-mean bound]\label{lem:step-009-small-mean}
Under Lemmas~\ref{lem:step-009-usage-law} and
\ref{lem:step-009-factorial-tail}, define

\[
\lambda:=\frac nk.
\tag{8}
\]

If \(\lambda\le2\), then the exact budget (1) satisfies \(M=8\) and

\[
\Pr\{U>M\}
\le\frac{\lambda^9}{9!}
\le\frac{2^9}{9!}
<2^{-9}.
\tag{9}
\]

If, more specifically, \(n<k\), then \(M=8\) and

\[
\Pr\{U>M\}<\frac1{9!}<2^{-9}.
\tag{10}
\]

Thus the statement includes the exact boundary \(\lambda=2\), every
occurrence of \(M=8\), and the case \(n<k\).
\end{lemma}

\begin{proof}
If \(\lambda\le2\), then \(4\lambda\le8\), so

\[
\lceil4n/k\rceil=\lceil4\lambda\rceil\le8.
\]

Equation (1) therefore gives \(M=8\) exactly. Since \(U\) is integer,

\[
\{U>M\}=\{U>8\}=\{U\ge9\}.
\tag{11}
\]

Apply Lemma~\ref{lem:step-009-factorial-tail} to the exact law in
Lemma~\ref{lem:step-009-usage-law} with \(p=1/k\) and \(r=9\). Because
\(np=n/k=\lambda\),

\[
\Pr\{U>M\}\le\frac{\lambda^9}{9!}
\le\frac{2^9}{9!}.
\tag{12}
\]

The final comparison in (9) is strict and completely numerical:

\[
9!=362880>262144=2^{18}.
\tag{13}
\]

Dividing (13) by the positive number \(2^9 9!\) gives
\(2^9/9!<2^{-9}\), completing (9). If \(n<9\), the event in (11) is
actually empty because \(U\le n\), and the convention in
Lemma~\ref{lem:step-009-factorial-tail} records this as \((n)_9=0\).

If \(n<k\), then \(0<\lambda<1\) and
\(0<4\lambda<4\), whence \(\lceil4\lambda\rceil\le4\) and \(M=8\).
The first inequality of (9) is then strict:

\[
\Pr\{U>M\}\le\frac{\lambda^9}{9!}<\frac1{9!}.
\]

Since \(9!=362880>512=2^9\), one has \(1/9!<2^{-9}\), proving (10).
Conversely, if \(M=8\), then (1) implies
\(\lceil4\lambda\rceil\le8\), hence \(4\lambda\le8\) and
\(\lambda\le2\); thus every possible \(M=8\) case is indeed contained
in this lemma. \end{proof}

\begin{lemma}[Monotone factorial envelope]\label{lem:step-009-envelope}
For every integer \(j\ge9\), define

\[
a_j:=\frac{(j/4)^{j+1}}{(j+1)!}.
\tag{14}
\]

Then \((a_j)_{j\ge9}\) is strictly decreasing and

\[
a_j\le a_9
=\frac{(9/4)^{10}}{10!}
<2^{-10}.
\tag{15}
\]
\end{lemma}

\begin{proof}
Direct cancellation gives, for every integer \(j\ge9\),

\[
\frac{a_{j+1}}{a_j}
=\frac{j+1}{4(j+2)}
 \left(1+\frac1j\right)^{j+1}.
\tag{16}
\]

We bound the power in (16) without importing an asymptotic exponential
estimate. The binomial theorem gives

\[
\left(1+\frac1j\right)^j
=\sum_{r=0}^j \binom jr j^{-r}
=\sum_{r=0}^j \frac{(j)_r}{r!\,j^r}
\le\sum_{r=0}^j\frac1{r!}
<3.
\tag{17}
\]

For the strict last inequality, \(r!\ge2^{r-1}\) for every \(r\ge2\),
with strict inequality for every \(r\ge3\), so

\[
\sum_{r=2}^{\infty}\frac1{r!}
<\sum_{r=2}^{\infty}\frac1{2^{r-1}}=1.
\]

Adding the \(r=0,1\) terms proves (17). Since \(j\ge9\),

\[
\left(1+\frac1j\right)^{j+1}
<3\left(1+\frac19\right)
=\frac{10}{3}<4.
\tag{18}
\]

Substituting (18) into (16) yields the strict ratio bound

\[
\frac{a_{j+1}}{a_j}
<\frac{j+1}{j+2}<1.
\tag{19}
\]

Thus the sequence is strictly decreasing and \(a_j\le a_9\). It remains
to verify the endpoint constant. Since \(4^{10}=2^{20}\), the inequality
\(a_9<2^{-10}\) is equivalent to

\[
9^{10}<2^{10}\,10!.
\]

Both integers are explicit:

\[
9^{10}=3486784401
<3715891200=2^{10}\,10!.
\tag{20}
\]

Equations (19)--(20) prove (15). \end{proof}

\begin{lemma}[Ceiling-controlled large-mean bound]\label{lem:step-009-large-mean}
Under Lemmas~\ref{lem:step-009-usage-law},
\ref{lem:step-009-factorial-tail}, and
\ref{lem:step-009-envelope}, if \(\lambda=n/k>2\) and

\[
j:=\lceil4\lambda\rceil,
\tag{21}
\]

then \(j\) is an integer satisfying \(j\ge9\), the exact budget is
\(M=j\), and

\[
\Pr\{U>M\}
\le\frac{\lambda^{j+1}}{(j+1)!}
\le\frac{(j/4)^{j+1}}{(j+1)!}
\le\frac{(9/4)^{10}}{10!}
<2^{-10}<2^{-9}.
\tag{22}
\]

The conclusion remains valid when \(j+1>n\), when the overflow event is
empty.
\end{lemma}

\begin{proof}
The strict inequality \(\lambda>2\) gives \(4\lambda>8\). Therefore its
integer ceiling obeys

\[
j=\lceil4\lambda\rceil\ge9,
\tag{23}
\]

so (1) gives \(M=\max\{8,j\}=j\). Because \(U\) is integer,

\[
\{U>M\}=\{U\ge j+1\}.
\tag{24}
\]

Apply Lemma~\ref{lem:step-009-factorial-tail} with
\(p=1/k\) and the integer order \(r=j+1\). Together with
Lemma~\ref{lem:step-009-usage-law}, this gives

\[
\Pr\{U>M\}
\le\frac{(n/k)^{j+1}}{(j+1)!}
=\frac{\lambda^{j+1}}{(j+1)!}.
\tag{25}
\]

If \(j+1>n\), the sharper first expression in (6) has
\((n)_{j+1}=0\), so (24) is empty and (25) still holds. No unstated
assumption \(j+1\le n\) is needed.

The ceiling definition (21) also gives \(j\ge4\lambda\), hence
\(\lambda\le j/4\). Raising nonnegative quantities to the positive integer
power \(j+1\) and dividing by \((j+1)!>0\) proves the second inequality in
(22). Lemma~\ref{lem:step-009-envelope} supplies the next two inequalities,
and \(2^{-10}<2^{-9}\) follows from \(2^{-10}=\tfrac12 2^{-9}\).
This proves the entire strict chain (22). \end{proof}

\begin{proposition}[Exact finite-budget overflow certificate]\label{prop:step-009-overflow}
Under the basic-setting ranges \(k\ge2,n\ge1\), the exact budget (1), and
the ideal hidden-coordinate tag interface, the usage count (2)
satisfies (3). More explicitly:

\begin{enumerate}
\item If \(n/k\le2\), then \(M=8\) and (9) gives the strict tail bound.
\item If \(n/k>2\), then \(M=\lceil4n/k\rceil\ge9\) and (22) gives the
stronger tail bound \(\Pr(U>M)<2^{-10}\).
\item If \(k=2\) or \(k=3\), then \(M\ge n\) and \(\Pr(U>M)=0\).
\item If \(n<k\), then \(M=8\) and the strict chain (10) holds.
\item If \(M=8\), necessarily \(n/k\le2\), so the first branch applies.
\end{enumerate}

No contradiction-regime, hard-prior, privacy, PAC, or learner condition is
needed for any clause.
\end{proposition}

\begin{proof}
Lemma~\ref{lem:step-009-usage-law} proves the first assertion in (3).
Exactly one of \(n/k\le2\) and \(n/k>2\) holds. In the first case,
Lemma~\ref{lem:step-009-small-mean} proves
\(\Pr(U>M)<2^{-9}\); in the second,
Lemma~\ref{lem:step-009-large-mean} proves the stronger
\(\Pr(U>M)<2^{-10}<2^{-9}\). This proves the second assertion in (3)
for all \(k,n\).

For the requested small-tag boundaries, if \(k\in\{2,3\}\), then
\(4/k\ge1\), so

\[
\left\lceil\frac{4n}{k}\right\rceil\ge n
\quad\text{and hence}\quad M\ge n.
\tag{26}
\]

The count (2) always obeys \(0\le U\le n\). Thus \(U>M\) is impossible
and its probability is zero. This direct argument covers every \(n\),
including values for which \(M=8\) and values for which the ceiling branch
exceeds eight.

The \(n<k\) and \(M=8\) clauses, including all ceiling inequalities and
their strict constants, were proved in
Lemma~\ref{lem:step-009-small-mean}. Finally, \(U=0\) never belongs to the
overflow event because \(M\ge8\); it remains an ordinary atom of the exact
binomial law rather than a conditioned-away case. \end{proof}
